% problem-000317 7 银席夫碱大环配合物
\section*{7. 银席夫碱大环配合物}
（图：）

\subsection*{7-1}
将纯化所得产物做了元素分析和红外光谱等基本表征，化合物的官能团在红外光谱上产⽣不同的特征吸收峰，如何利⽤配合物1的红外光谱判断产物中不含原料及低分⼦量的线性聚合物类的副产物？

\subsection*{7-2}
在氩⽓氛中得到了配合物 $^ { 3 , }$ ，将配合物1和3在氩⽓氛中⽤⾜量的 $\mathrm { N a B H _ { 4 } }$ 还原，得到化合物4和 $^ { 5 , }$ ，并进⾏了红外光谱研究。请推测化合物4和5红外光谱上最主要的差别是什么？

\subsection*{7-3}
测定了配合物1和3的晶体结构，结果发现配合物3中⽆氧，配合物1中和O原⼦相连的C原⼦的相关结构数据为：C—O 键⻓为 150.0 pm、C—N 键⻓为 124.7 pm，N—C—O、C—C—O 和 N—C—C 的键⻆分别为 $1 2 6 . 0 ^ { \circ }$ $1 1 2 . 5 ^ { \circ }$ 和 $1 2 1 . 5 ^ { \circ }$ ，如何利⽤这些数据说明在空⽓中得到的确实是配合物 1，⽽不是由于 $\mathrm { H _ { 2 } O }$ 分⼦在C=N双键上的加成所得到的配合物2？

\subsection*{7-4}
在盐酸存在下将配合物 1 和 3 ⽔解，过滤，滤液经分离除去 trpn 后的剩余组分经⾼效液相⾊谱分析，配合物1的⽔解产物⾊谱图上显⽰有A和B两个组分，代表纯组分相对量的峰⾯积⽐为33. 02:66.09 (A:B)，⽽配合物3的⽔解产物仅有⼀种组分 $\mathbf { C } _ { \circ }$ 。对A和B进⾏了质谱表征，其电喷雾质谱图分别为(a)和(b)（质谱图的纵坐标为相对丰度，横坐标为质荷⽐，相对丰度100%的峰对应的质荷⽐为�+1）。

（图：）
说明A、B和C各为何物。

（图：）

\subsection*{7-5}
在⽆ trpn、其他条件和合成配合物 1 相同时， $\mathrm { A g N O _ { 3 } }$ 在空⽓中不能氧化间苯⼆甲醛。对配合物进⾏氢核磁共振实验，所得谱图如下（图中不同位置的信号表⽰氢的化学环境不同），其中(a)图对应配合物1，(b)图对应⽆氧条件下新鲜制备的配合物3，(c)图对应配合物3在空⽓中放置⼀周后所得样品。结合前⾯⼏步的实验能否判断配合物1中氧原⼦的来源？请简述理由。

（图：）

（图：）

（图：）
